\scanpage{95}
\mspage{\unsure{88|89}}
% archive sheet number: 94 (upper left)
\deleted{$\unsure{J|I}$ de cardinal $\leq\pi$ et grand devant $\pi$, soit $I$}

J'ai envie de reprendre la définition de l'a\deleted{\unsure{cc}} $\pi$-accessibilité de $M$, pour une filtration cardinale donnée à coups de conditions équivalentes. Pour tt système inductif \deleted{\textcircled{$\unsure{a|\alpha}$}} $(X_i)_{i\in I}$ dans $M$, je garde la notation
\[
I_M=\{i\in I\mid X_i\in M\}
\]
Je pose, pour tt cardinal $\pi\geq\pi_0$
\[
M^{\pi}=M\cap\mathrm{Filt}^{\pi}
\]
\deleted{\textcircled{$\unsure{a_\pi|b}$} $M$}

\textcircled{$a_\pi$} $\unsure{M|\underline M}$ est l'ens des \deleted{\unsure{limites}} objets de la forme
\[
\varinjlim_{i\in I}X_i,
\]
où les $X_i$ sont dans $M^{\pi}$, et $I$ est grand devant $\pi$.

\textcircled{\textcircled{$b_\pi$}} Soit $(X_i)_{i\in I}$ un syst. inductif dans $M$, \deleted{$I$ grand devant $\pi$, et} \inserted{parfaitement adapté à la filtration cardinale, pour $\pi$, i.e. $I$ filtr., et toute partie filtrante $J$ de cardinal $\leq\pi$ admet une borne supérieure $i_J$, et $X_{i_J}\xleftarrow{\simeq}\varinjlim_J X_i$, et enfin $X_i\in\mathrm{Filt}^{\pi}$ pour tt $i\in I$.} Soit
\[
X=\varinjlim X_i.
\]
Pour que l'on ait $X\in M$, il f. et s. que $\unsure{I_M|I^{\pi}_M}$ soit cofinal dans $I$.

\textcircled{$c_\pi$} $\underline M^{\pi}=\underline M\cap\underline{\mathrm{Filt}}^{\pi}$ est une souscatégorie de $\underline M$ génératrice par épimorphismes stricts,

\textcircled{$d_\pi$} \deleted{[et par \unsure{limites}]} $M$ est stable par $\varinjlim$ grandes devant $\pi$, enfin pour tt $X\in M$, $\underline M^{\pi}/X$ est une catégorie filtrante.$^\ast$
% slanted
\margin{$\ast$ cf plus bas pour une \unsure{vérif.}}
